\section{Protocolo experimental y plan de estimación}\label{sec:protocolo}

Las secciones anteriores definieron el modelo (Sección~\ref{sec:modelo-general}),
el montaje (Sección~\ref{sec:diseno}) y el mapa de observación
(Sección~\ref{sec:mediciones}). Esta sección fija el experimento concreto:
la matriz de corridas, los valores numéricos de entradas y salidas, la
clasificación de los parámetros del modelo según su origen, y el esquema de
estimación basado en una red neuronal informada por física (PINN).

\subsection{Matriz experimental}\label{sec:matriz}

El experimento consta de cinco corridas secuenciales
(Observación~\ref{rem:escalado}), resumidas en la
Tabla~\ref{tab:matriz-experimental}. El diseño obedece a tres criterios de la
literatura de GEI con \textit{H.~illucens}: triplicado de la condición
principal \citep{lindbergProcessEfficiencyGreenhouse2022,pangInfluenceCarbonNitrogen2020},
lote testigo sin larvas para la separación microbiana
\citep{parodiBioconversionEfficienciesGreenhouse2020a,rossiEstimatingDynamicsGreenhouse2024a},
y una corrida de extrapolación que pone a prueba la capacidad predictiva del
modelo fuera de la condición de ajuste.

\begin{table}[htbp]
  \centering
  \caption{Matriz experimental. Todas las corridas usan el mismo sustrato,
  la misma dosis ($1{,}4$~g por larva) y el mismo calendario de cierres.}
  \label{tab:matriz-experimental}
  \begin{tabular}{clccccl}
    \hline
    Corrida & Rol & $N_0$ & $S_0$ (g) & $\theta_{s,0}$ & Duración & Réplica de \\
    \hline
    E0 & Caracterización (vacía) & 0 & 0 & --- & $\sim 2$~d & --- \\
    E1 & Testigo microbiano & 0 & 560 & $0{,}70$ & $t_f$ & --- \\
    E2 & Condición principal & 400 & 560 & $0{,}70$ & $t_f$ & --- \\
    E3 & Condición principal & 400 & 560 & $0{,}70$ & $t_f$ & E2 \\
    E4 & Condición principal & 400 & 560 & $0{,}70$ & $t_f$ & E2 \\
    E5 & Validación por extrapolación & 700 & 980 & $0{,}70$ & $t_f$ & --- \\
    \hline
  \end{tabular}
\end{table}

\begin{description}
  \item[E0: caracterización.] Cámara vacía; ejecuta las cuatro pruebas de la
  Sección~\ref{sec:caracterizacion} (hermeticidad, $V_{\mathrm{aire}}$,
  calibración, línea base de sensores). Se repite en la recalibración entre
  corridas.
  \item[E1: testigo.] Sustrato sin larvas ($N_0 = 0$). Mide
  $r^{\mathrm{mic}}_{\mathrm{CO_2}}$, $r^{\mathrm{mic}}_{\mathrm{O_2}}$ y
  $r_{\mathrm{CH_4}}$ en todo el horizonte, resolviendo la identificabilidad
  de la Sección~\ref{sec:identificabilidad}.
  \item[E2--E4: condición principal.] $N_0 = 400$, el escenario central de la
  Tabla~\ref{tab:carga-lote}. Las tres réplicas secuenciales estiman la
  variabilidad entre lotes de los parámetros ajustados.
  \item[E5: validación.] $N_0 = 700$, el escenario extremo. El modelo se
  ajusta con E2--E4 y debe \emph{predecir} E5 sin reajuste: es la prueba de
  generalización del gemelo digital.
\end{description}

\begin{remark}[Replicación secuencial y ruta de paralelización]\label{rem:paralelizacion}
Las réplicas son secuenciales por economía de equipamiento: una sola cámara,
un solo tren de sensores. El costo es el calendario ($\sim 4$~meses para las
cinco corridas) y la confusión corrida--tiempo, mitigada con metadatos por
corrida (cepa, edad y lote de neonatos), recalibración entre corridas y el
testigo E1 como referencia común. Si el proyecto dispone de recursos
adicionales, la paralelización no exige triplicar el tren de gases: basta
replicar la cámara y los sensores baratos (SHT35, sensor capacitivo, celda de
carga) y multiplexar el muestreo de gases con electroválvulas hacia el mismo
NDIR, esquema análogo al de las cuatro cámaras de
\citet{rossiEstimatingDynamicsGreenhouse2024a}.
\end{remark}

\subsection{Valores de entrada}\label{sec:entradas}

Las entradas del sistema son las consignas de la ley de control
(Definición~\ref{def41}), la ventilación conmutada (Definición~\ref{def42}) y
las condiciones iniciales del lote. La Tabla~\ref{tab:entradas} fija sus
valores numéricos.

\begin{table}[htbp]
  \centering
  \caption{Entradas del experimento: consignas, ventilación y condiciones
  iniciales (corridas E2--E4).}
  \label{tab:entradas}
  \begin{tabular}{llll}
    \hline
    Entrada & Símbolo & Valor & Fuente \\
    \hline
    Consigna de temperatura & $T_{\mathrm{ref}}$ & $28$~°C & zona de confort larval \\
    Consigna de humedad & $H_{\mathrm{ref}}$ & $70$~\% & operación del lecho \\
    Caudal operativo & $Q_{\mathrm{op}}$ & $1$~L/min & Ec.~\eqref{eq:q-dimensionamiento} \\
    Calendario de cierres & $\{t_k^c\}$ & cada $2$~días & Obs.~\ref{rem:cierres} \\
    Duración de cierre & $\Delta$ & Tabla~\ref{tab:delta-adaptativo} & Ecs.~\eqref{eq:delta-min}--\eqref{eq:delta-max} \\
    Carga de larvas & $N_0$ & $400$ neonatas & Tabla~\ref{tab:carga-lote} \\
    Carga de sustrato & $S_0$ & $560$~g ($1{,}4$~g/larva) & dosis del diseño \\
    Humedad inicial & $\theta_{s,0}$ & $0{,}70$ (gravimétrico) & determinación en carga \\
    Condiciones del inlet & $(T_{\mathrm{in}}, w_{\mathrm{in}}, c_{i,\mathrm{in}})$ & medidas & forzamientos conocidos \\
    \hline
  \end{tabular}
\end{table}

El criterio de terminación $t_f$ es el de la Ec.~\eqref{eq:tf}, evaluado en
línea por el gemelo digital mediante la firma metabólica
(Observación~\ref{rem:tf}), con confirmación visual de la fracción de
prepupas al abrir la cámara.

\subsection{Valores de salida}\label{sec:salidas}

La Tabla~\ref{tab:salidas} resume las mediciones con su frecuencia y rango
esperado. Las series continuas se registran a $0{,}1$~Hz; las muestras de
cromatografía se toman por el septum al inicio, mitad y fin de cada cierre
(Sección~\ref{sec:caracterizacion}).

\begin{table}[htbp]
  \centering
  \caption{Salidas medidas: variable, sensor, frecuencia y rango esperado.}
  \label{tab:salidas}
  \begin{tabular}{lllll}
    \hline
    Variable & Sensor & Frecuencia & Rango esperado & Estado asociado \\
    \hline
    $T$ aire & SHT35 & $0{,}1$~Hz & $20$--$35$~°C & $T$ \\
    $H$ aire & SHT35 & $0{,}1$~Hz & $40$--$95$~\% & $w$ \\
    $c_{\mathrm{CO_2}}$ & NDIR & $0{,}1$~Hz & $420$--$15\,000$~ppm & $c_{\mathrm{CO_2}}$ \\
    $c_{\mathrm{O_2}}$ & electroquímico & $0{,}1$~Hz & $15$--$21$~\% & $c_{\mathrm{O_2}}$ \\
    $c_{\mathrm{CH_4}}$ & NDIR & $0{,}1$~Hz & $0$--$500$~ppm & $c_{\mathrm{CH_4}}$ \\
    $(\theta_s^{\mathrm{sen}}, T_s^{\mathrm{sen}})$ & capacitivo & $0{,}1$~Hz & operación del lecho & $\theta_s$, $T_s$ \\
    $m_{\mathrm{LC}}$ & celda de carga & $0{,}1$~Hz & decreciente & restricción~\eqref{eq:celda-carga} \\
    CO$_2$, CH$_4$ & GC (septum) & $3$ por cierre & --- & validación metrológica \\
    H$_2$S, NH$_3$ & supervisión & $0{,}1$~Hz & umbrales de alarma & fuera del modelo \\
    \hline
  \end{tabular}
\end{table}

Los rangos de CO$_2$ y O$_2$ cubren el pico metabólico de la
Tabla~\ref{tab:delta-adaptativo}; el de CH$_4$ se fija bajo a propósito: en
condiciones aireadas la metanogénesis es marginal
\citep{rossiEstimatingDynamicsGreenhouse2024a}, y cualquier superación del
rango es en sí misma una señal de anaerobiosis
(Observación~\ref{rem:microbiano}).

\subsection{Clasificación de parámetros}\label{sec:clasificacion-parametros}

Los parámetros del modelo se clasifican en tres grupos según su origen
(Tabla~\ref{tab:parametros}). Esta clasificación es la base del plan de
estimación: solo el tercer grupo se ajusta con los datos del experimento.

\begin{table}[htbp]
  \centering
  \caption{Origen de los parámetros del modelo.}
  \label{tab:parametros}
  \begin{tabular}{lll}
    \hline
    Grupo & Parámetros & Fuente \\
    \hline
    Fijados & $a$, $m$, $Y_B$, $Y_L$, $B_{\max}$, $\beta$, $\eta_A$ &
    literatura \citep{eriksenDynamicModellingFeed2022,eriksenMetabolicPerformanceFeed2024} \\
    Fijados & $\gamma_q$, $\gamma_w$, $\lambda$ &
    equivalentes energéticos \citep{lightonMeasuringMetabolicRates2008,hansenUseCalorespirometricRatios2004} \\
    Medidos & $V_{\mathrm{aire}}$, $k_{\mathrm{fuga}}$, $h_\theta$, $\sigma_i$ &
    caracterización (Sección~\ref{sec:caracterizacion}) \\
    Estimados & $k_{\mathrm{mic}}$, $h_s A_s$, $C_{\mathrm{th}}$, $C_{\mathrm{th},s}$,
    coef.\ de $E_{\mathrm{lecho}}$ & ajuste con E1--E4 \\
    \hline
  \end{tabular}
\end{table}

\begin{remark}[Reajuste fino de parámetros fisiológicos]
Los parámetros fisiológicos se adoptan de la literatura porque fueron
ajustados sobre datos de crecimiento directos, que este experimento no
recolecta en línea. Si los residuos de validación (Sección~\ref{sec:estimacion})
resultan sistemáticos, se admite un reajuste fino de $(a, m)$ ---los dos
parámetros con mayor influencia sobre $r_{\mathrm{CO_2}}$ según la
Ec.~\eqref{eq:rco2-general}--- reportándolo explícitamente como desviación
del modelo de referencia.
\end{remark}

\subsection{Estimación con red neuronal informada por física}\label{sec:pinn}

La estimación del grupo de parámetros estimados y de los estados latentes
$(A, B, L, N, S)$ usa una PINN \citep{raissiPhysicsinformedNeuralNetworks2019a},
con la forma unificada ventilado--cerrado del modelo
(Sección~\ref{sec:modelo-general}) como restricción física.

\paragraph{Arquitectura y pérdida.} La red aproxima los estados como funciones
del tiempo, $\hat{\mathbf{x}}(t; \varphi)$, y la función de pérdida combina
tres términos:
\begin{equation}
  \mathcal{L} = \underbrace{\mathcal{L}_{\mathrm{datos}}}_{\text{ajuste a } \mathbf{y}(t)}
  + \underbrace{\mathcal{L}_{\mathrm{física}}}_{\text{residuos de las EDO}}
  + \underbrace{\mathcal{L}_{\mathrm{ci}}}_{\text{condiciones iniciales}},
  \label{eq:pinn-loss}
\end{equation}
donde $\mathcal{L}_{\mathrm{datos}}$ penaliza la discrepancia con las series
de la Tabla~\ref{tab:salidas} y con las tasas $R_i(t_k)$ del Nivel~1
(Sección~\ref{sec:estimacion}); $\mathcal{L}_{\mathrm{física}}$ evalúa los
residuos de cada balance del modelo en puntos de colocación dentro de
$[0, t_f]$; y $\mathcal{L}_{\mathrm{ci}}$ ancla las condiciones iniciales
conocidas (Tabla~\ref{tab:entradas}). Los parámetros estimados de la
Tabla~\ref{tab:parametros} son variables del entrenamiento junto con
$\varphi$.

\paragraph{Partición de datos.} Coherente con la
Observación~\ref{rem:id-lazo}: los cierres de ventilación y las ventanas de
consigna en escalón son los datos de \emph{identificación}; la operación
regulada normal es \emph{validación}. La corrida E5 queda reservada como
prueba de extrapolación final y no participa en ningún ajuste.

\paragraph{Referencia comparativa.} Un filtro de Kalman de ensamble sobre la
forma unificada sirve de línea base: la PINN debe al menos igualarlo en
residuos de validación para justificar su complejidad.

\paragraph{Criterios de éxito.} Los de la Sección~\ref{sec:estimacion}:
residuos compatibles con el ruido de sensores, cumplimiento de la
restricción de la celda de carga~\eqref{eq:celda-carga}, parámetros dentro
de rangos fisiológicos, y ---criterio adicional de esta sección--- error de
predicción de E5 del mismo orden que el error de validación de E2--E4.
